 % tkz-lib-eu-shape.tex
% Copyright 2011-2026  Alain Matthes
 % SPDX-License-Identifier: LPPL-1.3c
 % Maintainer: Alain Matthes

\typeout{2026/01/25 5.13c tkz-lib-eu-shape.tex}
%<--------------------------------------------------------------------------–>
%  Création des symboles
%<--------------------------------------------------------------------------–>
\makeatletter
% define a new shape for the points
\pgfdeclareshape{cross}
{%
  \inheritsavedanchors[from=rectangle] % this is nearly a rectangle
  \inheritanchorborder[from=rectangle]
  \inheritanchor[from=rectangle]{north}
  \inheritanchor[from=rectangle]{north west}
  \inheritanchor[from=rectangle]{north east}
  \inheritanchor[from=rectangle]{center}
  \inheritanchor[from=rectangle]{west}
  \inheritanchor[from=rectangle]{east}
  \inheritanchor[from=rectangle]{mid}
  \inheritanchor[from=rectangle]{mid west}
  \inheritanchor[from=rectangle]{mid east}
  \inheritanchor[from=rectangle]{base}
  \inheritanchor[from=rectangle]{base west}
  \inheritanchor[from=rectangle]{base east}
  \inheritanchor[from=rectangle]{south}
  \inheritanchor[from=rectangle]{south west}
  \inheritanchor[from=rectangle]{south east}
  \foregroundpath{
% store lower right in xa/ya and upper right in xb/yb
  \southwest \pgf@xa=\pgf@x \pgf@ya=\pgf@y
  \northeast \pgf@xb=\pgf@x \pgf@yb=\pgf@y
  \pgfpathmoveto{\pgfqpoint{0 pt}{\pgf@ya}}
  \pgfpathlineto{\pgfqpoint{0 pt}{\pgf@yb}}
  \pgfpathmoveto{\pgfqpoint{\pgf@xa}{0 pt}}
  \pgfpathlineto{\pgfqpoint{\pgf@xb}{0 pt}}
 }
}
\makeatother
\endinput